% Page budget: 1.55 pages | Owns: augmentation topology, additional states, encoder, and closed-loop objective.
% WRITING GUIDE
% Purpose: Define the final augmentation architecture and explain how it is trained inside the known feedback loop.
% Start here: Current implementation decisions D-129, D-140, D-141, D-180, D-220, D-233, and D-237 in docs/decisions.md.
% Supporting material: Quinten's 18 September outline feedback, the Meeting-audit synthesis, Thesis-writeup/Documentation/TRAINING-DESIGN.md, and docs/writeup figures.
% Mandatory papers: Read Hoekstra's 2025 LFR augmentation paper, the 2026 first-principles augmentation paper, the encoder-initialization paper, and Kessels' Chapter 5 before writing this section.
% Research-plan anchor: Preserve Aspect 2's dynamic parallel augmentation objective; document the departure from open-loop training and earlier state routing.
% Current decisions: Separate physical and additional states, keep the controller outside the model, score the complete training window without burn-in, and use one network for the physical correction and added-state update.
% Choices to justify: Final reported routing, added-state dimension, ANN width and depth, zero initialization, encoder history, closed-loop objective, rollout length, full-window scoring, and validation horizon. Use the configuration of the reported thesis run, not a superseded routing experiment.
% Horizon check: Treat encoder history, scored training window, joint-estimation horizon, and free-run validation horizon separately; relate candidates to relevant stable modes and test sensitivity to slower dynamics and free-integrator accumulation.
% Implementation audit: Read scripts/gantry/gantry_interconnect_dynamic.py, gantry_dynamic/{model,config,controller,data,training}.py, and model_augmentation/fit_systems/{interconnect,closed_loop,pre_encoder}.py.
% Reference check: Verify sources for parallel LFR augmentation, encoder initialization, closed-loop state extension, and any claimed architectural precedent against the original paper or thesis.
% Cautions: Earlier routing plans are superseded; distinguish the truth-only hidden absorber from learned states; closed-loop fit alone does not prove autonomous plant recovery.
% Figures: Absorber schematic, author's updated architecture, and thesis-owned common-reference loops. No residual panel or horizon figure.
% Section is finished when: Every signal has a defined frame and timing, the RK4 boundary is explicit, and the described architecture matches the final code.
%
% SOURCE MAP
% Hoekstra et al. 2025 EJC: dynamic-parallel topology, additional states, encoder-based truncated simulation.
% Hoekstra et al. 2026 Automatica submission: extended LFR formulation, baseline-preserving model/encoder initialization, joint identification context.
% Hoekstra et al. 2026 encoder paper: reconstructability-based W^b initialization; it does not identify W^a or test dynamic augmentation.
% Kessels 2025 thesis: direct closed-loop rollout, controller equations, and reconstruction of the machine controller state for each window.
% Drenth 2025 thesis: LPV-specific augmentation precedent; not the unrestricted MLP structure used here.
% This project: gantry specialization, residual-loop implementation, exact routing, RK4 boundary, and verification.
\section{Dynamic LPV-LFR Augmentation}\label{sec:augmentation}

The baseline of Section~\ref{sec:lfr} describes the rigid-body dynamics of
the gantry.
Garc\'ia-Herreros et al.\ neglect the resonance of the supporting base, the
vibration of the cross-arm on its flexible plates, the detent forces of the
linear motors, and the variation of friction along the
stroke~\cite[Sec.~2.4]{garcia2013model}.
The baseline \eqref{eq:eom} also omits the Coulomb friction of their
model~\cite{garcia2013model}.
We augment the baseline with learned dynamics.
The physical parameters $\theta_{\mathrm{base}}$, the network parameters
$\theta_{\mathrm{aug}}$ and the encoder parameters $\eta$ are estimated
jointly from records measured under feedback.
Section~\ref{sec:aug_structure} defines the augmented model,
Section~\ref{sec:aug_encoder} the initial state of each training window, and
Section~\ref{sec:aug_closed_loop} the closed-loop objective.

\subsection{Augmented model}\label{sec:aug_structure}
If the omitted effects are dynamic, the six baseline states cannot represent
them.
Refitting $\theta_{\mathrm{base}}$ adds no states, and neither does a static
correction of the state update~\cite[Sec.~2]{hoekstra2025lfr}.
The augmentation therefore needs states of its own.
We use the dynamic-parallel structure~\cite[Table~1]{hoekstra2026lfr}
\begin{equation}
\begin{aligned}
 \tilde x_{k+1}
 &=f_{\mathrm{base}}(\theta_{\mathrm{base}},\tilde x_k,\hat u_k)
   +f_{\mathrm{aug}}(\theta_{\mathrm{aug}},\hat x_k,\hat u_k),\\
 \bar x_{k+1}
 &=g_{\mathrm{aug}}(\theta_{\mathrm{aug}},\hat x_k,\hat u_k),\\
 \hat y_k&=C_n\tilde x_k ,
\end{aligned}
\label{eq:aug_transition}
\end{equation}
where $\tilde x=\col(q,\dot q)\in\R^6$ is the physical state of
Section~\ref{sec:lfr}, $\bar x\in\R^{n_{x_a}}$ collects the additional
states, and $\hat x=\col(\tilde x,\bar x)$.
The output $\hat y$ estimates the measured positions $y=\qact$, and $C_n$ is
their kinematic map $\begin{bmatrix}P^\top & 0\end{bmatrix}$ from
\eqref{eq:Ptransform}.
The input $\hat u=P\hat u_{\mathrm{act}}$ is the generalised force of the
actuator forces $\hat u_{\mathrm{act}}$ applied in the closed-loop rollout of
Section~\ref{sec:aug_closed_loop}.
Drenth also augments a self-scheduled LPV-LFR baseline with additional
states, with a learning component that is itself an
LPV-LFR~\cite[Sec.~5.2]{drenth2025thesis}.
Here, the learning component is a multilayer perceptron (MLP) of the
augmented state and the input (Fig.~\ref{fig:aug_structure}).

$f_{\mathrm{base}}$ is one fourth-order Runge-Kutta (RK4) step of
\eqref{eq:eom} over the sample period $T_s$, with the input held over the
step.
Each stage evaluates $M(Y)$ at its own state, so $f_{\mathrm{base}}$ remains
self-scheduled.
All signals are normalised to zero mean and unit variance with the
statistics of the training records, and $f_{\mathrm{base}}$ and $C_n$ act in
these coordinates~\cite[Sec.~5.3]{hoekstra2026lfr}.

The two learned functions are the outputs of one MLP with tanh hidden
layers and a linear output layer,
\begin{equation}
 \begin{bmatrix}
  f_{\mathrm{aug}}(\theta_{\mathrm{aug}},\hat x_k,\hat u_k)\\
  g_{\mathrm{aug}}(\theta_{\mathrm{aug}},\hat x_k,\hat u_k)
 \end{bmatrix}
 =\phi_{\mathrm{aug}}(\theta_{\mathrm{aug}},z_k),\qquad
 z_k=\col(\hat x_k,\hat u_{\mathrm{act},k}),
 \label{eq:aug_mlp}
\end{equation}
where $\theta_{\mathrm{aug}}$ collects the weights and biases, with $W_L$
and $b_L$ those of the output layer.
The first six outputs form $f_{\mathrm{aug}}$.
As in~\cite[Sec.~2]{hoekstra2025lfr}, the learning component acts on the
entire state, the positions included, so it can also correct the $X$ and
$Y$ rows into which omitted effects may couple.
The last $n_{x_a}$ outputs form $g_{\mathrm{aug}}$ and give the next
additional state.
Because the network acts on the result of the RK4 step, $f_{\mathrm{aug}}$
is a discrete-time correction, whereas $f_{\mathrm{base}}$ integrates the
continuous-time physics.
The output map is not learned, because the measured positions follow
exactly from the physical state~\cite[Remark~5.3]{kessels2025ai}.

\begin{figure}[t]
 \centering
 \resizebox{\columnwidth}{!}{\input{figures/tex/augmentation_structure}}
 \caption{Augmented model. The physics step $f_{\mathrm{base}}$ and the
 network $\phi_{\mathrm{aug}}$ act in parallel on $\hat x_k$.
 $\phi_{\mathrm{aug}}$ supplies $f_{\mathrm{aug}}$ and $g_{\mathrm{aug}}$,
 and the output map $h_{\mathrm{base}}=C_n$ is not learned
 ($h_{\mathrm{aug}}=0$). The encoder $\psi$ sets $\hat x$ at each window
 start. During training, $u_k$ is the closed-loop input $\hat u_k$.}
 \label{fig:aug_structure}
\end{figure}

Within one step, no learned function feeds the input of another, so the
computational graph of \eqref{eq:aug_transition} is acyclic.
$f_{\mathrm{base}}$ is smooth because $M(Y)$ stays nonsingular for every $Y$
(Section~\ref{sec:params}).
With tanh activations, the model is therefore well
posed~\cite[Thm.~9]{hoekstra2026lfr}.

The additional states reach the output only through $f_{\mathrm{aug}}$, one
step later.
As in any state-space model, $\bar x$ is determined only up to an invertible
state transformation.
The additional states therefore have no assigned physical meaning, also when
the omitted effect is a physical mode.

Training starts from the baseline.
The output layer of the network is initialised at zero,
\begin{equation}
 W_L=0,\quad b_L=0
 \quad\Longrightarrow\quad
 f_{\mathrm{aug}}=0,\quad g_{\mathrm{aug}}=0 ,
 \label{eq:aug_zero_init}
\end{equation}
so the augmented model reproduces the baseline from the same physical
initial state~\cite[Eq.~(29)]{hoekstra2026lfr}.
From this start, $f_{\mathrm{aug}}$ can also learn changes that a different
$\theta_{\mathrm{base}}$ would produce, so the split between the two
components is not unique~\cite[Sec.~5.2]{hoekstra2026lfr}.
Section~\ref{sec:obc} constrains this split.

\subsection{Encoder and initial state}\label{sec:aug_encoder}
Each training window simulates \eqref{eq:aug_transition} from an initial
state, but only the positions are measured.
Following Beintema et al.~\cite{beintema2023subnet}, an encoder estimates
this state from the recorded input and output history,
\begin{equation}
 \hat x_{\tau|\tau}
 =\psi_\eta(\xi_\tau),\qquad
 \xi_\tau=\col\bigl(y_{\tau-n:\tau},\,u_{\mathrm{act},\tau-n:\tau}\bigr),
 \label{eq:aug_encoder}
\end{equation}
where $\tau$ is the window start, $n$ the encoder lag,
$u_{\mathrm{act}}$ the recorded actuator force, and $y_{\tau-n:\tau}$
stacks the samples from $\tau-n$ to $\tau$.
The history ends at the current sample, which the reconstructability map
below uses~\cite[Eq.~(15)]{hoekstra2026encoder}.
The encoder estimates all states, the measured positions included.

The encoder is a linear map with a nonlinear
correction~\cite[Eq.~(8)]{hoekstra2026encoder},
\begin{equation}
 \psi_\eta(\xi_\tau)
 =\begin{bmatrix}W^b\\ W^a\end{bmatrix}\xi_\tau+\tilde\psi(\xi_\tau),
 \label{eq:aug_enc_lin}
\end{equation}
where $W^b$ and $W^a$ give the physical and additional states and
$\tilde\psi$ is an MLP.

The physical rows $W^b$ are initialised from the baseline.
Linearising \eqref{eq:eom} at the nominal parameters about rest gives
\begin{equation}
 A_c(Y)=\begin{bmatrix}0 & I\\ -M(Y)^{-1}K & -M(Y)^{-1}C\end{bmatrix},\qquad
 B_c(Y)=\begin{bmatrix}0\\ M(Y)^{-1}P\end{bmatrix},
 \label{eq:aug_lin}
\end{equation}
with state $\tilde x$ and input $u_{\mathrm{act}}$.
With $(A_d,B_d)$ the zero-order-hold discretisation of $(A_c(0),B_c(0))$
over $T_s$ and $C_d=\begin{bmatrix}P^\top & 0\end{bmatrix}$, $W^b$ is the
reconstructability map~\cite[Eqs.~(16), (17)]{hoekstra2026encoder}
\begin{equation}
 W^b=\begin{bmatrix}
  A_d^nO_n^{+} & \;r_n-A_d^nO_n^{+}T_n
 \end{bmatrix},
 \label{eq:aug_wb}
\end{equation}
where $O_n$ is the observability matrix of $(A_d,C_d)$ over the encoder lag
$n$, $T_n$ the Toeplitz matrix of input-to-output responses, $r_n$ the
input-to-state map over the same lag, and the columns follow the output and
input parts of $\xi_\tau$.
All matrices are normalised as the signals.
The measured positions make the linear model observable, so $O_n$ has full
column rank and $O_n^{+}$ is a left inverse.
The linearisation at $Y=0$ only sets the starting value of $W^b$, because
the encoder is trained jointly with the model.

The rows $W^a$ have no baseline counterpart and are drawn at random, and
the correction $\tilde\psi$ starts at zero.
By \eqref{eq:aug_zero_init}, the encoded additional states do not affect the
output at initialisation.
They become active as the output layer of the network is trained.

\subsection{Closed-loop rollout and objective}\label{sec:aug_closed_loop}
The records are measured under feedback.
The model is used to predict the machine under feedback: the response to a given reference, feedforward and
known controller.
We therefore fit this closed-loop response, rather than an open-loop
simulation driven by the recorded input as
in~\cite[Eq.~(22)]{hoekstra2026lfr}.

The baseline gives a second reason.
$X$ and $Y$ have no stiffness in \eqref{eq:damping_stiffness_matrices}, so
$A_c(Y)$ in \eqref{eq:aug_lin} has two eigenvalues at the origin for every
$Y$.
In an open-loop rollout, a velocity error therefore integrates into a
position error that persists.
Inside the loop, the controller acts on such errors as it does on the
machine, provided the model in feedback with the controller is stable.
The controller attenuates genuine model error in the same way, so a small
closed-loop error does not show that the model is accurate in open loop.

Kessels trains an augmented model with the known feedback controller inside
each truncated rollout and propagates gradients through
it~\cite[Eqs.~(5.12), (5.13c), (5.13d)]{kessels2025ai}.
We adopt this and keep the controller outside the learned model, so that it
can be replaced in evaluation~\cite[Sec.~5.3.2.3]{kessels2025ai}.

The recorded loop is driven by a reference $r$ and a feedforward
$u^{\mathrm{ff}}$ through a linear time-invariant controller with state $s$.
The model loop receives the same $r$, $u^{\mathrm{ff}}$ and controller, and
differs only in its output $\hat y$ (Fig.~\ref{fig:aug_closed_loop}):
\begin{equation}
\begin{aligned}
 s_{k+1}&=A_{\mathrm{fb}} s_k+B_{\mathrm{fb}}(r_k-y_k), \\
 u_{\mathrm{act},k}&=C_{\mathrm{fb}} s_k+D_{\mathrm{fb}}(r_k-y_k)+u^{\mathrm{ff}}_k,\\
 \hat s_{k+1}&=A_{\mathrm{fb}} \hat s_k+B_{\mathrm{fb}}(r_k-\hat y_k), \\
 \hat u_{\mathrm{act},k}&=C_{\mathrm{fb}} \hat s_k+D_{\mathrm{fb}}(r_k-\hat y_k)+u^{\mathrm{ff}}_k ,
\end{aligned}
\label{eq:aug_direct_controller}
\end{equation}
where $(A_{\mathrm{fb}},B_{\mathrm{fb}},C_{\mathrm{fb}},D_{\mathrm{fb}})$
realises the controller.
If both loops apply the same controller and feedforward and no actuator
saturates, subtracting the recorded loop from the model loop removes $r$ and
$u^{\mathrm{ff}}$,
\begin{equation}
\begin{aligned}
 x^c_{k+1}&=A_{\mathrm{fb}} x^c_k+B_{\mathrm{fb}}(y_k-\hat y_k),\qquad x^c_\tau=0,\\
 \hat u_{\mathrm{act},k}&=u_{\mathrm{act},k}+C_{\mathrm{fb}} x^c_k+D_{\mathrm{fb}}(y_k-\hat y_k),
\end{aligned}
\label{eq:aug_controller}
\end{equation}
where $x^c=\hat s-s$ is the difference of the controller states.
The model thus receives the recorded input plus the controller's response to
its own output error.

\begin{figure}[t]
 \centering
 \includegraphics[width=\columnwidth]{closed_loop_same_reference.pdf}
 \caption{Recorded plant loop (top) and augmented-model loop (bottom),
 driven by one shared reference $r_k$ and feedforward $u_k^{\mathrm{ff}}$,
 with the same feedback controller. The model output is evaluated
 before the input advances its state to the next sample.}
 \label{fig:aug_closed_loop}
\end{figure}

The condition $x^c_\tau=0$ starts the model's controller in the state of the
recorded controller, $\hat s_\tau=s_\tau$, which the recorded input already
contains.
Model error before $\tau$ is therefore not carried into the window.
Kessels instead reconstructs the controller state from the measured output,
the reference and the controller~\cite[Remark~5.4]{kessels2025ai}.
The form \eqref{eq:aug_controller} needs no such reconstruction.

Because the model has no feedthrough, each sample is evaluated in a fixed
order: $\hat y_k$ from the state, then $\hat u_{\mathrm{act},k}$ from
\eqref{eq:aug_controller}, then both states advance.

Over window starts $\mathcal T$ and window length $n_f$, the parameters
$\vartheta=(\theta_{\mathrm{base}},\theta_{\mathrm{aug}},\eta)$ minimise the
normalised output error
\begin{equation}
 V(\vartheta)=\frac{1}{|\mathcal T|\,n_f}
 \sum_{\tau\in\mathcal T}\sum_{\ell=0}^{n_f-1}
 \norm{y_{\tau+\ell}-\hat y_{\tau+\ell|\tau}}_2^2
 \label{eq:aug_loss}
\end{equation}
subject to \eqref{eq:aug_transition}, \eqref{eq:aug_encoder} and
\eqref{eq:aug_controller}.
This is the truncated objective of~\cite[Eq.~(22)]{hoekstra2026lfr}, with
each window simulated in closed loop as
in~\cite[Eq.~(5.12)]{kessels2025ai}.
As in both, every sample of the window is scored.
The physical parameters, the network and the encoder are estimated jointly,
while the controller is fixed.
Section~\ref{sec:setup} gives $n_{x_a}$, the network size, the encoder lag
$n$, the window length $n_f$ and the horizon over which the models are
validated.
